% ═══════════════════════════════════════════════════════════
% §1 Introduction
% ═══════════════════════════════════════════════════════════
\section{Introduction}
\label{sec:intro}

Residue--residue contact prediction from multiple sequence alignments has
transformed protein structure determination~\cite{weigt2009,morcos2011,
	jones2012,kamisetty2013}.  The dominant paradigm, Direct Coupling Analysis
(DCA), models the joint distribution of amino-acid sequences through a Potts
Hamiltonian whose coupling parameters are estimated by inverting a covariance
matrix computed from deep alignments~\cite{morcos2011}.  Sparse inverse
covariance estimation, as implemented in PSICOV~\cite{jones2012}, achieves a
mean top-$L/5$ precision of approximately 0.73 on standard benchmarks.
Pseudo-likelihood maximisation~\cite{ekeberg2013} and related methods further
improve upon mean-field approximations.  More recently, end-to-end deep
learning approaches such as AlphaFold2~\cite{jumper2021} have largely
superseded standalone contact prediction by integrating coevolutionary signal
with structural priors inside attention-based neural architectures.

Despite this success, the encoding layer through which biological sequences
enter the statistical machinery remains opaque.  No prior work has
established machine-checkable proofs that the transformation from amino-acid
alphabet to the numerical representation used for inference preserves the
structural properties that downstream analysis requires.  This gap is not
merely academic: subtle encoding errors can silently corrupt results while
producing plausible-looking output.

Here we address this gap by constructing a formally verified bridge between
digital-logic design theory and protein coevolution analysis.  Our central
observation is that the Karnaugh map~\cite{karnaugh1953}, when combined with
reflected binary Gray-code ordering of amino acids~\cite{gray1953}, provides
a natural embedding of aligned sequences into a grid structure where adjacent
cells correspond to physicochemically similar residue pairs.  The
Quine--McCluskey algorithm~\cite{quine1952,mccluskey1956} then extracts
minimal sum-of-products expressions from thresholded frequency maps,
producing prime implicants that we interpret as coevolutionary constraints.

We formalise 118 theorems in Lean~4 covering Gray-code bijectivity, $k$-mer
indexing injectivity, amino-acid encoding properties, contact-map topology,
implicant-cover soundness, and padding safety.  We apply the verified pipeline
to 150 protein domains from the PSICOV benchmark comprising over 1.6 million
labelled observations against native experimental structures.  Our results
demonstrate that Boolean circuits transfer across proteins, that their
cell-rate prior improves DCA ranking when fused via rank aggregation, and
that negative-selection (forbidden-combination) circuits generalise with high
specificity.  We also report negative results honestly: an attention-style
tension mechanism does not improve prediction, iterative direct information
is unreliable on conserved-rich families, and current circuit precision falls
far short of the level required for template-free structure reconstruction.

All software, formal proofs, result files, and documentation are publicly
available (\texttt{github.com/E-Motioner-X-SBS/co-evolution-analysis}).
